\section{Suites numériques}

\subsection{Définitions et modes de génération}

\begin{definition}
Une suite numérique $(u_n)_{n \in \mathbb{N}}$ est une application de $\mathbb{N}$ (ou une partie de $\mathbb{N}$) dans $\mathbb{R}$.
\end{definition}

Une suite peut être définie de manière :
\begin{itemize}
\item \textbf{Explicite :} $u_n = f(n)$ pour tout $n \in \mathbb{N}$.
\item \textbf{Récurrente :} $u_0$ donné et $u_{n+1} = f(u_n)$ pour tout $n \in \mathbb{N}$.
\end{itemize}

\subsection{Suites arithmétiques}

\begin{definition}
Une suite $(u_n)$ est arithmétique de raison $r$ si pour tout $n \in \mathbb{N}$ :
$$u_{n+1} = u_n + r$$
\end{definition}

\textbf{Terme général :} $u_n = u_0 + nr$

\textbf{Somme des $n$ premiers termes :}
$$S_n = \sum_{k=0}^{n} u_k = (n+1) \cdot \dfrac{u_0 + u_n}{2}$$

\subsection{Suites géométriques}

\begin{definition}
Une suite $(u_n)$ est géométrique de raison $q$ ($q \neq 0$) si pour tout $n \in \mathbb{N}$ :
$$u_{n+1} = q \cdot u_n$$
\end{definition}

\textbf{Terme général :} $u_n = u_0 \cdot q^n$

\textbf{Somme des $n$ premiers termes :}
Si $q \neq 1$ :
$$S_n = u_0 \cdot \dfrac{1 - q^{n+1}}{1 - q}$$

\subsection{Convergence}

\begin{theoreme}
Toute suite croissante et majorée converge.\\
Toute suite décroissante et minorée converge.
\end{theoreme}

\begin{theoreme}[Suites adjacentes]
Si $(u_n)$ est croissante, $(v_n)$ est décroissante, et $\lim_{n \to +\infty}(v_n - u_n) = 0$,
alors $(u_n)$ et $(v_n)$ convergent vers la même limite.
\end{theoreme}

\subsection{Limites usuelles}

\begin{itemize}
\item $\lim_{n \to +\infty} q^n = 0$ si $|q| < 1$
\item $\lim_{n \to +\infty} \dfrac{n^k}{a^n} = 0$ pour tout $k \in \mathbb{N}$ et $a > 1$
\item $\lim_{n \to +\infty} \dfrac{\ln n}{n} = 0$
\end{itemize}
